\ifdefined\gkver\else\def\gkver{student}\fi
\documentclass[\gkver]{gaokaozhenti}
\newcommand{\ccnum}[1]{{\fontspec{Noto Sans CJK JP}#1}}
\begin{document}

\chapter{2015年山东理综}
%% paperType: 理综物理部分

\begin{enumerate}
\item
%% number: 14
%% paperName: 2015年山东理综第 14 题 6 分
%% typeId: 1
%% type: 单选题
%% chapter: 曲线运动
%% point: 平抛运动
%% method: 无
%% score: 6
%% degree: 650
%% duplicateId: 0
%% body:
距地面高 $5\Um$ 的水平直轨道 A、B 两点相距 $2\Um$，在 B 点用细线悬挂一小球，离地高度为 $h$，如图。小车始终以 $4\Ums$ 的速度沿轨道匀速运动，经过 A 点时将随车携带的小球由轨道高度自由卸下，小车运动至 B 点时细线被轧断，最后两球同时落地。不计空气阻力，取重力加速度的大小 $g=10\Umsq$。可求得 $h$ 等于\xzanswer{A}

\onepicture[h=2.6cm]{figs/14.png}

\fourchoices[answer=A]
{$1.25\Um$}
{$2.25\Um$}
{$3.75\Um$}
{$4.75\Um$}

%% answer: A
%% memo:
\memoanswer{
	小车上的物体落地的时间 $t=\sqrt{\frac{2H}{g}}$；小车从 A 到 B 的时间 $t_{1}=\frac{x}{v}$，小球下落的时间 $t_{2}=\sqrt{\frac{2h}{g}}$。由题意得时间关系：$t=t_{1}+t_{2}$，即 $\sqrt{\frac{2H}{g}}=\frac{d}{v}+\sqrt{\frac{2h}{g}}$，解得：$h=1.25\Um$，选项 A 正确。
}

\item
%% number: 15
%% paperName: 2015年山东理综第 15 题 6 分
%% typeId: 1
%% type: 单选题
%% chapter: 曲线运动
%% point: 人造地球卫星
%% method: 无
%% score: 6
%% degree: 650
%% duplicateId: 0
%% body:
如图，拉格朗日点 $L_{1}$ 位于地球和月球连线上，处在该点的物体在地球和月球引力的共同作用下，可与月球一起以相同的周期绕地球运动。据此，科学家设想在拉格朗日点 $L_{1}$ 建立空间站，使其与月球同周期绕地球运动。以 $a_{1}$、$a_{2}$ 分别表示该空间站和月球向心加速度的大小，$a_{3}$ 表示地球同步卫星向心加速度的大小。以下判断正确的是\xzanswer{D}

\onepicture[h=2.4cm]{figs/15.png}

\fourchoices[answer=D]
{$a_{2}>a_{3}>a_{1}$}
{$a_{2}>a_{1}>a_{3}$}
{$a_{3}>a_{1}>a_{2}$}
{$a_{3}>a_{2}>a_{1}$}

%% answer: D
%% memo:
\memoanswer{
	由于 $G\frac{Mm}{r^{2}}=ma_{n}$，则 $a_{n}=\frac{GM}{r^{2}}$，由题意可知月球的轨道半径大于空间站的轨道半径，而同步卫星的运动周期小于空间站的周期，故空间站的轨道半径大于同步卫星的轨道半径，所以 $a_{3}>a_{2}>a_{1}$。选项 D 正确。
}

\item
%% number: 16
%% paperName: 2015年山东理综第 16 题 6 分
%% typeId: 1
%% type: 单选题
%% chapter: 运动和力
%% point: 连接体
%% method: 无
%% score: 6
%% degree: 450
%% duplicateId: 0
%% body:
如图，滑块 A 置于水平地面上，滑块 B 在一水平力作用下紧靠滑块 A（A、B 接触面竖直），此时 A 恰好不滑动，B 刚好不下滑。已知 A 与 B 间的动摩擦因数为 $\mu_{1}$，A 与地面间的动摩擦因数为 $\mu_{2}$，最大静摩擦力等于滑动摩擦力。A 与 B 的质量之比为\xzanswer{B}

\onepicture[h=2.4cm]{figs/16.png}

\fourchoices[answer=B]
{$\frac{1}{\mu_{1}\mu_{2}}$}
{$\frac{1-\mu_{1}\mu_{2}}{\mu_{1}\mu_{2}}$}
{$\frac{1+\mu_{1}\mu_{2}}{\mu_{1}\mu_{2}}$}
{$\frac{2+\mu_{1}\mu_{2}}{\mu_{1}\mu_{2}}$}

%% answer: B
%% memo:
\memoanswer{
	物体 AB 整体在水平方向受力平衡，则有 $F=\mu_{2}(m_{A}+m_{B})g$；对于物体 B 在竖直方向平衡有：$\mu_{1}F=m_{B}g$，联立解得：$\frac{m_{A}}{m_{B}}=\frac{1-\mu_{1}\mu_{2}}{\mu_{1}\mu_{2}}$，选项 B 正确。
}

\item
%% number: 17
%% paperName: 2015年山东理综第 17 题 6 分
%% typeId: 2
%% type: 多选题
%% chapter: 电磁感应
%% point: 右手定则
%% method: 无
%% score: 6
%% degree: 450
%% duplicateId: 0
%% body:
如图，一均匀金属圆盘绕通过其圆心且与盘面垂直的轴逆时针匀速转动。现施加一垂直穿过圆盘的有界匀强磁场，圆盘开始减速。在圆盘减速过程中，以下说法正确的是\xzanswer{ABD}

\onepicture[h=2.6cm]{figs/17.png}

\fourchoices[answer=ABD]
{处于磁场中的圆盘部分，靠近圆心处电势高}
{所加磁场越强越易使圆盘停止转动}
{若所加磁场反向，圆盘将加速转动}
{若所加磁场穿过整个圆盘，圆盘将匀速转动}

%% answer: ABD
%% memo:
\memoanswer{
	由右手定则可知，处于磁场中的圆盘部分，靠近圆心处电势高，选项 A 正确；根据 $E=BLv$ 可知所加磁场越强，则感应电动势越大，感应电流越大，产生的电功率越大，消耗的机械能越快，则圆盘越容易停止转动，选项 B 正确；若加反向磁场，根据楞次定律可知安培力阻碍圆盘的转动，故圆盘仍减速转动，选项 C 错误；若所加磁场穿过整个圆盘则圆盘中无感应电流，不消耗机械能，圆盘匀速转动。选项 D 正确，选项 ABD 正确。
}

\item
%% number: 18
%% paperName: 2015年山东理综第 18 题 6 分
%% typeId: 1
%% type: 单选题
%% chapter: 电场
%% point: 电场叠加
%% method: 无
%% score: 6
%% degree: 650
%% duplicateId: 0
%% body:
直角坐标系 $xOy$ 中，M、N 两点位于 $x$ 轴上，G、H 两点坐标如图，M、N 两点各固定一负点电荷，一电量为 $Q$ 的正点电荷置于 $O$ 点时，G 点处的电场强度恰好为零。静电力常量用 $k$ 表示。若将该正点电荷移到 G 点，则 H 点处场强的大小和方向分别为\xzanswer{B}

\onepicture[h=3.2cm]{figs/18.png}

\fourchoices[answer=B]
{$\frac{3kQ}{4a^{2}}$，沿 $y$ 轴正向}
{$\frac{3kQ}{4a^{2}}$，沿 $y$ 轴负向}
{$\frac{5kQ}{4a^{2}}$，沿 $y$ 轴正向}
{$\frac{5kQ}{4a^{2}}$，沿 $y$ 轴负向}

%% answer: B
%% memo:
\memoanswer{
	因正电荷在 O 点时，G 点场强为零，则两负电荷在 G 点形成的电场的合场强为 $E_{1}=k\frac{Q}{a^{2}}$；若将正电荷移到 G 点，则正电荷在 H 点的场强为 $E_{2}=k\frac{Q}{(2a)^{2}}=\frac{1}{4}\frac{kQ}{a^{2}}$，因两负电荷在 G 点的场强与在 H 点的场强等大反向，则 H 点的合场强为 $E=E_{1}-E_{2}=\frac{3kQ}{4a^{2}}$，方向沿 $y$ 轴负向，选项 B 正确。
}

\item
%% number: 19
%% paperName: 2015年山东理综第 19 题 6 分
%% typeId: 1
%% type: 单选题
%% chapter: 电磁感应
%% point: 电磁感应中图象
%% method: 无
%% score: 6
%% degree: 450
%% duplicateId: 0
%% body:
如图甲，$R_{0}$ 为定值电阻，两金属圆环固定在同一绝缘平面内。左端连接在一周期为 $T_{0}$ 的正弦交流电源上，经二极管整流后，通过 $R_{0}$ 的电流 $i$ 始终向左，其大小按图乙所示规律变化。规定内圆环 a 端电势高于 b 端时，a、b 间的电压为 $u_{ab}$ 正，下列 $u_{ab}-t$ 图像可能正确的是\xzanswer{C}

\twopicture[showlabel=false, hA=2.4cm, hB=2.4cm]
{figs/19a.png}
{figs/19b.png}

\fourchoices[answer=C, ispicture=true, h=2.0cm]
{figs/19c1.png}
{figs/19c2.png}
{figs/19c3.png}
{figs/19c4.png}

%% answer: C
%% memo:
\memoanswer{
	在第一个 $0.25T_{0}$ 时间内，通过大圆环的电流为瞬时针逐渐增加，由楞次定律可判断内环内 a 端电势高于 b 端，而电流的变化率逐渐减小，所以内环的电动势逐渐减小；同理在第 $0.25T_{0}\sim0.5T_{0}$ 时间内，通过大圆环的电流为瞬时针逐渐减小，由楞次定律可判断内环内 a 端电势低于 b 端，而电流的变化率逐渐变大，所以内环的电动势逐渐变大，由题所给选项图判定，选项 C 正确。
}

\item
%% number: 20
%% paperName: 2015年山东理综第 20 题 6 分
%% typeId: 2
%% type: 多选题
%% chapter: 电场
%% point: 带电粒子的偏转
%% method: 无
%% score: 6
%% degree: 450
%% duplicateId: 0
%% body:
如图甲，两水平金属板间距为 $d$，板间电场强度的变化规律如图乙所示。$t=0$ 时刻，质量为 $m$ 的带电微粒以初速度 $v_{0}$ 沿中线射入两板间，$0\sim T/3$ 时间内微粒匀速运动，$T$ 时刻微粒恰好经金属边缘飞出。微粒运动过程中未与金属板接触。重力加速度的大小为 $g$。关于微粒在 $0\sim T$ 时间内运动的描述，正确的是\xzanswer{BC}

\twopicture[showlabel=false, hA=2.4cm, hB=2.8cm]
{figs/20a.png}
{figs/20b.png}

\fourchoices[answer=BC]
{末速度大小为 $\sqrt{2}v_{0}$}
{末速度沿水平方向}
{重力势能减少了 $mgd/2$}
{克服电场力做功为 $mgd$}

%% answer: BC
%% memo:
\memoanswer{
	因 $0\sim T/3$ 内微粒匀速运动，则 $qE_{0}=mg$；在 $T/3\sim2T/3$ 时间内，微粒只受重力作用，做平抛运动，在 $2T/3$ 时刻的竖直分速度 $v_{y1}=gT/3$，水平分速度为 $v_{0}$；在 $2T/3\sim T$ 时间内有 $2qE_{0}-mg=ma$，得 $a=g$，方向向上，在 $T$ 时刻，微粒的竖直分速度减小到 0，水平分速度仍为 $v_{0}$，A 错误，B 正确；微粒的重力势能减小了 $\Delta E_{p}=mgd/2$，C 正确；从射入到离开，由动能定理得：$mgd/2+W_{\text{电}}=0$，克服电场力做功 $mgd/2$，D 错误，选项 BC 正确。
}

\item
%% number: 21
%% paperName: 2015年山东理综第 21 题 10 分
%% typeId: 4
%% type: 实验题
%% chapter: 力物体的平衡
%% point: 验证平行四边形实验
%% method: 无
%% score: 10
%% degree: 450
%% duplicateId: 0
%% body:
某同学通过下述实验验证力的平行四边形定则。

实验步骤：

①将弹簧秤固定在贴有白纸的竖直木板上，使其轴线沿竖直方向。

②如图甲所示，将环形橡皮筋一端挂在弹簧秤的秤钩上，另一端用圆珠笔尖竖直向下拉，直到弹簧秤示数为某一设定值时，将橡皮筋两端的位置记为 $O_{1}$、$O_{2}$，记录弹簧秤的示数 $F$，测量并记录 $O_{1}$、$O_{2}$ 间的距离（即橡皮筋的长度 $l$）。每次将弹簧秤示数改变 $0.50\UN$，测出所对应的 $l$，部分数据如下表所示：

\begin{center}
\begin{tabular}{|c|c|c|c|c|c|c|}
\hline
$F$（$\UN$） & 0 & 0.50 & 1.00 & 1.05 & 2.00 & 2.50 \\
\hline
$l$（$\Ucm$） & $l_{0}$ & 10.97 & 12.02 & 13.00 & 13.98 & 15.05 \\
\hline
\end{tabular}
\end{center}

③找出②中 $F=2.50\UN$ 时橡皮筋两端的位置，重新记为 O、$O'$，橡皮筋的拉力记为 $F_{OO'}$。

④在秤钩上涂抹少许润滑油，将橡皮筋搭在秤钩上，如图乙所示。用两圆珠笔尖成适当角度同时拉橡皮筋的两端，使秤钩的下端达到 O 点，将两笔尖的位置标记为 A、B，橡皮筋 OA 段的拉力记为 $F_{OA}$，OB 段的拉力记为 $F_{OB}$。

完成下列作图和填空：

（1）利用表中数据在给出的坐标纸上（见答题卡）画出 $F-l$ 图线，根据图线求得 $l_{0}=$\tkanswer[1.4cm]{10.00}$\Ucm$。

（2）测得 $OA=6.00\Ucm$，$OB=7.60\Ucm$，则 $F_{OA}$ 的大小为\tkanswer[1.4cm]{1.80}$\UN$。

（3）根据给出的标度，在答题卡上作出 $F_{OA}$ 和 $F_{OB}$ 的合力 $F'$ 的图示。

（4）通过比较 $F'$ 与\tkanswer[1.6cm]{$F_{OO'}$}的大小和方向，即可得出实验结论。

\twopicture[showlabel=false, hA=4.0cm, hB=4.0cm]
{figs/21a.png}
{figs/21b.png}

%% answer:
\jdanswer{
（1）$10.00\Ucm$（9.8、9.9、10.1 均正确）；（2）$1.80\UN$（$1.70\sim1.90$ 均正确）；（3）如图所示；（4）$F_{OO'}$。

\onepicture[h=3.4cm]{figs/21c.png}
}
%% memo:
\memoanswer{
	（1）做出 $F-l$ 图像，求得直线的截距即为 $l_{0}$，可得 $l_{0}=10.00\Ucm$；

	\onepicture[h=3.0cm]{figs/21d.png}

	（2）可计算弹簧的劲度系数为 $k=\frac{F}{\Delta x}=\frac{1}{0.02}\UNm=50\UNm$；若 $OA=6.00\Ucm$，$OB=7.60\Ucm$，则弹簧的弹力 $F=k\Delta l=50\times(6.00+7.60-10.00)\times10^{-2}\UN=1.8\UN$；则此时 $F_{OA}=F_{OB}=F=1.8\UN$；

	（3）如图；

	（4）通过比较 $F'$ 和 $F_{OO'}$ 的大小和方向，可得出实验结论。
}

\item
%% number: 22
%% paperName: 2015年山东理综第 22 题 8 分
%% typeId: 4
%% type: 实验题
%% chapter: 电路
%% point: 输出功率极值
%% method: 无
%% score: 8
%% degree: 450
%% duplicateId: 0
%% body:
如图甲所示的电路图中，恒流源可作为电路提供恒定电流 $I_{0}$，$R$ 为定值电阻，电流表、电压表均可视为理想电表。某同学利用该电路研究滑动变阻器 $R_{L}$ 消耗的电功率。改变 $R_{L}$ 的阻值，记录多组电流、电压的数值，得到如图乙所示的 $U-I$ 关系图线。

回答下列问题：

（1）滑动触头向下滑动时，电压表示数将\tkanswer[1.2cm]{减小}（填“增大”或“减小”）；

（2）$I_{0}=$\tkanswer[1.2cm]{1.0}$\UA$；

（3）$R_{L}$ 消耗的最大功率\tkanswer[1.6cm]{5}$\UW$（保留一位有效数字）。

\twopicture[showlabel=false, hA=2.6cm, hB=2.4cm]
{figs/22a.png}
{figs/22b.png}

%% answer:
\jdanswer{
（1）减小；（2）$1.0\UA$；（3）$5\UW$。
}
%% memo:
\memoanswer{
	（1）滑动头向下移动时，$R_{L}$ 阻值减小，则总电阻减小，电压变小，则电压表读数变小；

	（2）由电路图可知：$I_{0}=I+\frac{U}{R}$，即：$U=I_{0}R-IR$，由 $U-I$ 图线可知，$I_{0}R=20$；而曲线的斜率为 $R$ 的大小，$R=k=\frac{20}{1.0}\UO=20\UO$，则 $I_{0}=1.0\UA$；

	（3）$R_{L}$ 消耗的功率为 $P=IU=20I-20I^{2}$，当 $I=0.5\UA$ 时，功率的最大值为 $P_{m}=5\UW$。
}

\item
%% number: 23
%% paperName: 2015年山东理综第 23 题 16 分
%% typeId: 6
%% type: 计算题
%% chapter: 曲线运动
%% point: 竖直平面圆周运动
%% method: 无
%% score: 16
%% degree: 450
%% duplicateId: 0
%% body:
如图甲所示，物块与质量为 $m$ 的小球通过不可伸长的轻质细绳跨过两等高定滑轮连接。物块置于左侧滑轮正下方的表面水平的压力传感装置上，小球与右侧滑轮的距离为 $l$。开始时物块和小球均静止，将此时传感装置的示数记为初始值。现给小球施加一始终垂直于 $l$ 段细绳的力，将小球缓慢拉起至细绳与竖直方向成 $60^\circ$ 角，如图乙所示，此时传感装置的示数为初始值的 1.25 倍；再将小球由静止释放，当运动至最低位置时，传感装置的示数为初始值的 0.6 倍。不计滑轮的大小和摩擦，重力加速度的大小为 $g$。求：

\begin{enumerate}
	\item
	物块的质量；

	\item
	从释放到运动至最低位置的过程中，小球克服阻力所做的功。
\end{enumerate}

\onepicture[align=right, h=2.6cm]{figs/23a.png}

\onepicture[align=right, h=2.6cm]{figs/23b.png}

%% answer:
\jdanswer{
\begin{enumerate}
	\item
	$3m$；

	\item
	$0.1mgl$。
\end{enumerate}
}
%% memo:
\memoanswer{
	（1）设开始时细绳的拉力大小为 $T_{1}$，传感装置的初始值为 $F_{1}$，物块质量为 $M$，由平衡条件得

	对小球，$T_{1}=mg$\ccnum{①}

	对物块，$F_{1}+T_{1}=Mg$\ccnum{②}

	当细绳与竖直方向的夹角为 $60^\circ$ 时，设细绳的拉力大小为 $T_{2}$，传感装置的示数为 $F_{2}$，据题意可知，$F_{2}=1.25F_{1}$，由平衡条件得

	对小球，$T_{2}=mg\cos60^\circ$\ccnum{③}

	对物块，$F_{2}+T_{2}=Mg$\ccnum{④}

	联立\ccnum{①②③④}式，代入数据得 $M=3m$\ccnum{⑤}

	（2）设小球运动至最低位置时速度的大小为 $v$，从释放到运动至最低位置的过程中，小球克服阻力所做的功为 $W_{t}$，由动能定理得

	$mgl(1-\cos60^\circ)-W_{t}=\frac{1}{2}mv^{2}$\ccnum{⑥}

	在最低位置，设细绳的拉力大小为 $T_{3}$，传感装置的示数为 $F_{3}$，据题意可知，$F_{3}=0.6F_{1}$，对小球，由牛顿第二定律得

	$T_{3}-mg=m\frac{v^{2}}{l}$\ccnum{⑦}

	对物块，由平衡条件得 $F_{3}+T_{3}=Mg$\ccnum{⑧}

	联立\ccnum{①②⑤⑥⑦⑧}式，代入数据得

	$W_{t}=0.1mgl$\ccnum{⑨}
}

\item
%% number: 24
%% paperName: 2015年山东理综第 24 题 20 分
%% typeId: 6
%% type: 计算题
%% chapter: 磁场
%% point: 带电粒子在磁场中的运动
%% method: 无
%% score: 20
%% degree: 250
%% duplicateId: 0
%% body:
如图所示，直径分别为 $D$ 和 $2D$ 的同心圆处于同一竖直面内，O 为圆心，GH 为大圆的水平直径。两圆之间的环形区域（Ⅰ区）和小圆内部（Ⅱ区）均存在垂直圆面向里的匀强磁场。间距为 $d$ 的两平行金属极板间有一匀强电场，上极板开有一小孔。一质量为 $m$，电量为 $+q$ 的粒子由小孔下方 $d/2$ 处静止释放，加速后粒子以竖直向上的速度 $v$ 射出电场，由 H 点紧靠大圆内侧射入磁场。不计粒子的重力。

\begin{enumerate}
	\item
	求极板间电场强度的大小；

	\item
	若粒子运动轨迹与小圆相切，求Ⅰ区磁感应强度的大小；

	\item
	若Ⅰ区，Ⅱ区磁感应强度的大小分别为 $2mv/qD$，$4mv/qD$，粒子运动一段时间后再次经过 H 点，求这段时间粒子运动的路程。
\end{enumerate}

\onepicture[align=right, h=3.4cm]{figs/24.png}

%% answer:
\jdanswer{
\begin{enumerate}
	\item
	$\frac{mv^{2}}{qd}$

	\item
	$\frac{4mv}{qD}$ 或 $\frac{4mv}{3qD}$

	\item
	$5.5\pi D$
\end{enumerate}
}
%% memo:
\memoanswer{
	（1）设极板间电场强度的大小为 $E$，对粒子在电场中的加速运动，由动能定理得

	$qE\frac{d}{2}=\frac{1}{2}mv^{2}$\ccnum{①}

	由\ccnum{①}式得 $E=\frac{mv^{2}}{qd}$\ccnum{②}

	（2）设Ⅰ区磁感应强度的大小为 $B$，粒子做圆周运动的半径为 $R$，由牛顿第二定律得

	$qvB=m\frac{v^{2}}{R}$\ccnum{③}

	如图甲所示，粒子运动轨迹与小圆相切有两种情况。若粒子轨迹与小圆外切，由几何关系得

	$R=\frac{D}{4}$\ccnum{④}

	联立\ccnum{③④}式得 $B=\frac{4mv}{qD}$\ccnum{⑤}

	\onepicture[h=3.0cm]{figs/24a.png}

	若粒子轨迹与小圆内切，由几何关系得

	$R=\frac{3D}{4}$\ccnum{⑥}

	联立\ccnum{③⑥}式得 $B=\frac{4mv}{3qD}$\ccnum{⑦}

	（3）设粒子在Ⅰ区和Ⅱ区做圆周运动的半径分别为 $R_{1}$、$R_{2}$，由题意可知Ⅰ和Ⅱ区磁感应强度的大小分别为 $B_{1}=\frac{2mv}{qD}$、$B_{2}=\frac{4mv}{qD}$，由牛顿第二定律得

	$qvB_{1}=m\frac{v^{2}}{R_{1}}$，$qvB_{2}=m\frac{v^{2}}{R_{2}}$\ccnum{⑧}

	代入数据得 $R_{1}=\frac{D}{2}$，$R_{2}=\frac{D}{4}$\ccnum{⑨}

	设粒子在Ⅰ区和Ⅱ区做圆周运动的周期分别为 $T_{1}$、$T_{2}$，由运动学公式得

	$T_{1}=\frac{2\pi R_{1}}{v}$，$T_{2}=\frac{2\pi R_{2}}{v}$\ccnum{⑩}

	据题意分析，粒子两次与大圆相切的时间间隔内，运动轨迹如图乙所示，根据对称可知，Ⅰ区两段圆弧所对圆心角相同，设为 $\theta_{1}$，Ⅱ区内圆弧所对圆心角设为 $\theta_{2}$，圆弧和大圆的两个切点与圆心 O 连线间的夹角设为 $\alpha$，有几何关系得

	$\theta_{1}=120^\circ$\ccnum{⑪}

	$\theta_{2}=180^\circ$\ccnum{⑫}

	$\alpha=60^\circ$\ccnum{⑬}

	\onepicture[h=3.0cm]{figs/24b.png}

	粒子重复上述交替运动回到 H 点，轨迹如图丙所示，设粒子在Ⅰ区和Ⅱ区做圆周运动的时间分别为 $t_{1}$、$t_{2}$，可得

	$t_{1}=\frac{360^\circ}{\alpha}\times\frac{\theta_{1}\times2}{360^\circ}T_{1}$，$t_{2}=\frac{360^\circ}{\alpha}\times\frac{\theta_{2}}{360^\circ}T_{2}$\ccnum{⑭}

	设粒子运动的路程为 $s$，由运动学公式得

	$s=v(t_{1}+t_{2})$\ccnum{⑮}

	\onepicture[h=3.0cm]{figs/24c.png}

	联立\ccnum{⑨⑩⑪⑫⑬⑭⑮}式得 $s=5.5\pi D$\ccnum{⑯}
}

\item
%% number: 37
%% paperName: 2015年山东理综第 37 题 8 分
%% typeId: 6
%% type: 计算题
%% chapter: 气体性质
%% point: 查理定律
%% method: 无
%% score: 8
%% degree: 650
%% duplicateId: 0
%% body:
【物理——选修 3-3】（12 分）

（1）（4 分）墨滴入水，扩而散之，徐徐混匀。关于该现象的分析正确的是\xzanswer{BC}。（双选，填正确答案标号）

\fourchoices[answer=BC]
{混合均匀主要是由于碳粒受重力作用}
{混合均匀的过程中，水分子和碳粒都做无规则运动}
{使用碳粒更小的墨汁，混合均匀的过程进行得更迅速}
{墨汁的扩散运动是由于碳粒和水分子发生化学反应引起的}

（2）（8 分）扣在水平桌面上的热杯盖有时会发生被顶起的现象。如图，截面积为 $S$ 的热杯盖扣在水平桌面上，开始时内部封闭气体的温度为 $300\UK$，压强为大气压强 $p_{0}$。当封闭气体温度上升至 $303\UK$ 时，杯盖恰好被整体顶起，放出少许气体后又落回桌面，其内部压强立即减为 $p_{0}$，温度仍为 $303\UK$。再经过一段时间，内部气体温度恢复到 $300\UK$。整个过程中封闭气体均可视为理想气体。求：

\begin{enumerate}
	\item
	当温度上升到 $303\UK$ 且尚未放气时，封闭气体的压强；

	\item
	当温度恢复到 $300\UK$ 时，竖直向上提起杯盖所需的最小力。
\end{enumerate}

\onepicture[align=right, h=1.8cm]{figs/37.png}

%% answer:
\jdanswer{
\begin{enumerate}
	\item
	$\frac{101}{100}p_{0}$

	\item
	$\frac{201}{10100}p_{0}S$
\end{enumerate}
}
%% memo:
\memoanswer{
	（1）根据分子动理论的知识可知，混合均匀主要是由于水分子做无规则运动，碳粒的无规则运动为布朗运动；由于布朗运动的剧烈程度与颗粒大小和温度有关，所以使用碳粒更小的墨汁，布朗运动会越明显，则混合均匀的过程进行得更迅速，选项 B、C 正确。

	（2）（i）以开始封闭的气体为研究对象，由题意可知，初状态温度 $T_{0}=300\UK$，压强为 $p_{0}$，末状态温度 $T_{1}=303\UK$，压强设为 $p_{1}$，由查理定律得

	$\frac{p_{0}}{T_{0}}=\frac{p_{1}}{T_{1}}$\ccnum{①}

	代入数据得 $p_{1}=\frac{101}{100}p_{0}$\ccnum{②}

	（ii）设杯盖的质量为 $m$，刚好被顶起时，由平衡条件得

	$p_{1}S=p_{0}S+mg$\ccnum{③}

	放出少许气体后，以杯盖内的剩余气体为研究对象，由题意可知，初状态温度 $T_{2}=303\UK$，压强 $p_{2}=p_{0}$，末状态温度 $T_{3}=300\UK$，压强设为 $p_{3}$，由查理定律得

	$\frac{p_{2}}{T_{2}}=\frac{p_{3}}{T_{3}}$\ccnum{④}

	设提起杯盖所需的最小力为 $F$，由平衡条件得

	$F+p_{3}S=p_{0}S+mg$\ccnum{⑤}

	联立\ccnum{②③④⑤}式，代入数据得 $F=\frac{201}{10100}p_{0}S$\ccnum{⑥}
}

\item
%% number: 38
%% paperName: 2015年山东理综第 38 题 8 分
%% typeId: 6
%% type: 计算题
%% chapter: 光学
%% point: 光的折射
%% method: 无
%% score: 8
%% degree: 450
%% duplicateId: 0
%% body:
【物理——选修 3-4】（12 分）

（1）（4 分）如图，轻弹簧上端固定，下端连接一小物块，物块沿竖直方向做简谐运动。以竖直向上为正方向，物块简谐运动的表达式为 $y=0.1\sin(2.5\pi t)\Um$。$t=0$ 时刻，一小球从距物块 $h$ 高处自由落下；$t=0.6\Us$ 时，小球恰好与物块处于同一高度。取重力加速度的大小为 $g=10\Umsq$。以下判断正确的是\xzanswer{AB}（双选，填正确答案标号）

\onepicture[h=3.8cm]{\includesvg{figs/38a}}

\fourchoices[answer=AB]
{$h=1.7\Um$}
{简谐运动的周期是 $0.8\Us$}
{$0.6\Us$ 内物块运动的路程是 $0.2\Um$}
{$t=0.4\Us$ 时，物块与小球运动方向相反}

（2）（8 分）半径为 $R$、介质折射率为 $n$ 的透明圆柱体，过其轴线 $OO'$ 的截面如图所示。位于截面所在平面内的一细束光线，以角 $i_{0}$ 由 O 点入射，折射光线由上边界的 A 点射出。当光线在 O 点的入射角减小至某一值时，折射光线在上边界的 B 点恰好发生全反射。求 A、B 两点间的距离。

\onepicture[align=right, h=2.6cm]{figs/38b.png}

%% answer:
\jdanswer{
$\left(\frac{1}{\sqrt{n^{2}-1}}-\frac{\sqrt{n^{2}-\sin^{2}i_{0}}}{\sin i_{0}}\right)R$
}
%% memo:
\memoanswer{
	（1）$t=0.6\Us$ 时，物块的位移为 $y=0.1\sin(2.5\pi\times0.6)\Um=-0.1\Um$；而小球 $h+|y|=\frac{1}{2}gt^{2}$，解得 $h=1.7\Um$，A 正确；简谐运动周期为 $T=\frac{2\pi}{\omega}=\frac{2\pi}{2.5\pi}\Us=0.8\Us$，B 正确；$0.6\Us$ 内物块运动的路程为 $3A=0.3\Um$，C 错误；$t=0.4\Us=\frac{T}{2}$，此时物块在平衡位置向下振动与小球运动方向相同，D 错误，选项 A、B 正确。

	（2）当光线在 O 点的入射角为 $i_{0}$ 时，设折射角为 $\gamma_{0}$，由折射定律得

	$\frac{\sin i_{0}}{\sin\gamma_{0}}=n$\ccnum{①}

	设 A 点与左端面的距离为 $d_{A}$，由几何关系得

	\onepicture[h=3.0cm]{figs/38c.png}

	$\sin\gamma_{0}=\frac{R}{\sqrt{d_{A}^{2}+R^{2}}}$\ccnum{②}

	若折射光线恰好发生全反射，则在 B 点的入射角恰好为临界角 C，设 B 点与左端面的距离为 $d_{B}$，由折射定律得

	$\sin C=\frac{1}{n}$\ccnum{③}

	由几何关系得 $\sin C=\frac{d_{B}}{\sqrt{d_{B}^{2}+R^{2}}}$\ccnum{④}

	设 A、B 两点间的距离为 $d$，可得 $d=d_{B}-d_{A}$\ccnum{⑤}

	联立\ccnum{①②③④⑤}式得 $d=\left(\frac{1}{\sqrt{n^{2}-1}}-\frac{\sqrt{n^{2}-\sin^{2}i_{0}}}{\sin i_{0}}\right)R$\ccnum{⑥}
}

\item
%% number: 39
%% paperName: 2015年山东理综第 39 题 8 分
%% typeId: 6
%% type: 计算题
%% chapter: 运动和力
%% point: 动量和能量
%% method: 无
%% score: 8
%% degree: 450
%% duplicateId: 0
%% body:
【物理——选修 3-5】（12 分）

（1）（4 分）$\ce{^{14}C}$ 发生放射性衰变为 $\ce{^{14}N}$，半衰期约为 5700 年。已知植物存活其间，其体内 $\ce{^{14}C}$ 与 $\ce{^{12}C}$ 的比例不变；生命活动结束后，$\ce{^{14}C}$ 的比例持续减少。现通过测量得知，某古木样品中 $\ce{^{14}C}$ 的比例正好是现代植物所制样品的二分之一。下列说法正确的是\xzanswer{AC}（双选，填正确答案标号）

\fourchoices[answer=AC]
{该古木的年代距今约为 5700 年}
{$\ce{^{12}C}$、$\ce{^{13}C}$、$\ce{^{14}C}$ 具有相同的中子数}
{$\ce{^{14}C}$ 衰变为 $\ce{^{14}N}$ 的过程中放出 $\beta$ 射线}
{增加样品测量环境的压强将加速 $\ce{^{14}C}$ 的衰变}

（2）（8 分）如图，三个质量相同的滑块 A、B、C，间隔相等地静置于同一水平轨道上。现给滑块 A 向右的初速度 $v_{0}$，一段时间后 A 与 B 发生碰撞，碰后 AB 分别以 $v_{0}/8$、$3v_{0}/4$ 的速度向右运动，B 再与 C 发生碰撞，碰后 B、C 粘在一起向右运动。滑块 A、B 与轨道间的动摩擦因数为同一恒定值。两次碰撞时间极短。求 B、C 碰后瞬间共同速度的大小。

\onepicture[align=right, h=1.6cm]{figs/39.png}

%% answer:
\jdanswer{
$v=\frac{\sqrt{21}}{16}v_{0}$
}
%% memo:
\memoanswer{
	（1）古木样品中 $\ce{^{14}C}$ 的比例正好是现代植物所制样品的二分之一，则可知经过的时间为一个半衰期，即该古木的年代距今约为 5700 年，A 正确；$\ce{^{12}C}$、$\ce{^{13}C}$、$\ce{^{14}C}$ 具有相同的质子数，由于质量数不同，故中子数不同，B 错误；根据核反应方程可知，$\ce{^{14}C}$ 衰变为 $\ce{^{14}N}$ 的过程中放出电子，即发出 $\beta$ 射线，C 正确；外界环境不影响放射性元素的半衰期，D 错误；选项 A、C 正确。

	（2）设滑块质量为 $m$，A 与 B 碰撞前 A 的速度为 $v_{A}$，由题意知，碰撞后 A 的速度 $v_{A}'=\frac{1}{8}v_{0}$，B 的速度 $v_{B}=\frac{3}{4}v_{0}$，由动量守恒定律得

	$mv_{A}=mv_{A}'+mv_{B}$\ccnum{①}

	设碰撞前 A 克服轨道阻力所做的功为 $W_{A}$，由功能关系得

	$W_{A}=\frac{1}{2}mv_{0}^{2}-\frac{1}{2}mv_{A}^{2}$\ccnum{②}

	设 B 与 C 碰撞前 B 的速度为 $v_{B}'$，B 克服轨道阻力所做的功为 $W_{B}$，由功能关系得

	$W_{B}=\frac{1}{2}mv_{B}^{2}-\frac{1}{2}mv_{B}'^{2}$\ccnum{③}

	据题意可知 $W_{A}=W_{B}$\ccnum{④}

	设 B、C 碰后瞬间共同速度的大小为 $v$，由动量守恒定律得

	$mv_{B}'=2mv$\ccnum{⑤}

	联立\ccnum{①②③④⑤}式，代入数据得 $v=\frac{\sqrt{21}}{16}v_{0}$\ccnum{⑥}
}

\end{enumerate}

\end{document}
